\documentclass[../../../include/open-logic-section]{subfiles}

\begin{document}

\olfileid{sth}{ordinals}{ordtype}
\olsection{الأعداد الترتيبية بوصفها أنماط ترتيب}

وقد ضممنا الاستبدال إلى مسلّماتنا، فنحن الآن في~$\ZFminus$، وبوسعنا
أن نقيم البرهان الذي قصدنا إليه.

\begin{thm}\ollabel{thmOrdinalRepresentation}
لكل ترتيب جيد عدد ترتيبي وحيد يتماثل معه.
\end{thm}

\begin{proof}
ليكن~$\tuple{\OLId{latin-looped}{A}, <}$ ترتيبًا جيدًا. تدل~\olref[basic]{ordisoidentity}
على أنه لا يتماثل مع أكثر من عدد ترتيبي واحد. ولنفرض، طلبًا للخلف،
أن~$\tuple{\OLId{latin-looped}{A}, <}$ لا يتماثل مع \emph{أي} عدد ترتيبي. أول ما نفعله
أن نصغّر~$\tuple{\OLId{latin-looped}{A}, <}$ قدر المستطاع. فإن وجدت قطعة ابتدائية حقيقية
$\tuple{\OLId{latin-looped}{A}_\OLId{latin-ordinary}{a}, <_\OLId{latin-ordinary}{a}}$ لا تتماثل مع عدد ترتيبي، أخذنا أصغر~$\OLId{latin-ordinary}{a} \in \OLId{latin-looped}{A}$
لهذه الخاصية، ووضعنا~$\OLId{latin-looped}{B} = \OLId{latin-looped}{A}_\OLId{latin-ordinary}{a}$ و$\mathord{\lessdot} = \mathord{<_\OLId{latin-ordinary}{a}}$.
وإن لم توجد، وضعنا~$\OLId{latin-looped}{B} = \OLId{latin-looped}{A}$ و$\mathord{\lessdot} = \mathord{<}$.

في كلا الحالين تتماثل كل قطعة ابتدائية حقيقية من~$\OLId{latin-looped}{B}$ مع عدد ترتيبي،
وهو وحيد لما تقدم. ومن ثم يضمن الاستبدال وجود المجموعة الآتية، وهي دالة:
\[
\OLId{latin-ordinary}{f} = \Setabs{\tuple{\beta, \OLId{latin-ordinary}{b}}}{\OLId{latin-ordinary}{b} \in \OLId{latin-looped}{B}\text{ و }
\ordeq{\beta}{\tuple{\OLId{latin-looped}{B}_\OLId{latin-ordinary}{b}, \lessdot_\OLId{latin-ordinary}{b}}}}
\]
ويتم الخلف إذا أثبتنا أن~$\OLId{latin-ordinary}{f}$ تماثل~$\OLId{greek-variable}{alpha} \to \OLId{latin-looped}{B}$ لعدد ترتيبي~$\OLId{greek-variable}{alpha}$.

فأما~$\ran{\OLId{latin-ordinary}{f}} = \OLId{latin-looped}{B}$ فواضح. وأما~$\OLId{latin-ordinary}{f}$ فتحفظ الترتيب بمقتضى
\olref[iso]{lemordsegments}؛ أي إن~$\gamma \in \beta$ إذا وفقط إذا
كان~$\OLId{latin-ordinary}{f}(\gamma) \lessdot \OLId{latin-ordinary}{f}(\beta)$. وبقي أن~$\dom{\OLId{latin-ordinary}{f}}$ عدد ترتيبي؛
وبحسب~\olref[basic]{corordtransitiveord} يكفينا إثبات تعدي~$\dom{\OLId{latin-ordinary}{f}}$.
خذ إذن~$\beta \in \dom{\OLId{latin-ordinary}{f}}$، أي~$\ordeq{\beta}{\tuple{\OLId{latin-looped}{B}_\OLId{latin-ordinary}{b}, \lessdot_\OLId{latin-ordinary}{b}}}$
لبعض~$\OLId{latin-ordinary}{b}$. فإذا كان~$\gamma \in \beta$ كان~$\gamma \in \dom{\OLId{latin-ordinary}{f}}$،
بمقتضى~\olref[iso]{wellordinitialsegment}. وبالتعميم نحصل على
$\beta \subseteq \dom{\OLId{latin-ordinary}{f}}$.
\end{proof}

لقد أمكن الآن وضع التعريف الذي كنا نرومه منذ~\olref[vn]{sec}.

\begin{defn}
إذا كان~$\tuple{\OLId{latin-looped}{A}, <}$ ترتيبًا جيدًا، سمينا~$\ordtype{\OLId{latin-looped}{A}, <}$ نمط
ترتيبه؛ وهو العدد الترتيبي الوحيد~$\OLId{greek-variable}{alpha}$ الذي يحقق
$\ordeq{\tuple{\OLId{latin-looped}{A}, < }}{\OLId{greek-variable}{alpha}}$.
\end{defn}

ويقتضي هذا التعريف المبدأين المطلوبين:

\begin{cor}\ollabel{ordtypesworklikeyouwant}
إذا كان~$\tuple{\OLId{latin-looped}{A}, <}$ و$\tuple{\OLId{latin-looped}{B}, \lessdot}$ ترتيبين جيدين، فإن:
\begin{align*}
\ordtype{\OLId{latin-looped}{A}, <} = \ordtype{\OLId{latin-looped}{B}, \lessdot}&\text{ إذا وفقط إذا كان }
\ordeq{\tuple{\OLId{latin-looped}{A}, <}}{\tuple{\OLId{latin-looped}{B}, \lessdot}}\\
\ordtype{\OLId{latin-looped}{A}, <} \in \ordtype{\OLId{latin-looped}{B}, \lessdot}&\text{ إذا وفقط إذا كان }
\ordeq{\tuple{\OLId{latin-looped}{A}, <}}{\tuple{\OLId{latin-looped}{B}_\OLId{latin-ordinary}{b}, \lessdot_\OLId{latin-ordinary}{b}}}
\text{ لبعض }\OLId{latin-ordinary}{b} \in \OLId{latin-looped}{B}
\end{align*}
\end{cor}

\begin{proof}
أما التكافؤ الأول فمن~\olref[basic]{ordisoidentity}. وللثاني نضع
$\ordtype{\OLId{latin-looped}{A}, <} = \OLId{greek-variable}{alpha}$ و$\ordtype{\OLId{latin-looped}{B}, \lessdot} = \beta$،
ونأخذ~$\OLId{latin-ordinary}{f} \colon \beta \to \OLId{latin-looped}{B}$ الدالة التي يقوم عليها تماثل الترتيب
مع~$\tuple {\OLId{latin-looped}{B}, \lessdot}$. حينئذ:
\begin{align*}
\OLId{greek-variable}{alpha} \in \beta&\text{ إذا وفقط إذا كان }
\funrestrictionto{\OLId{latin-ordinary}{f}}{\OLId{greek-variable}{alpha}} \colon \OLId{greek-variable}{alpha} \to \OLId{latin-looped}{B}_{\OLId{latin-ordinary}{f}(\OLId{greek-variable}{alpha})}
\text{ تماثلًا}\\
&\text{ إذا وفقط إذا كان }
\ordeq{\tuple{\OLId{latin-looped}{A}, <}}{\tuple{\OLId{latin-looped}{B}_{\OLId{latin-ordinary}{f}(\OLId{greek-variable}{alpha})}, \lessdot_{\OLId{latin-ordinary}{f}(\OLId{greek-variable}{alpha})}}}\\
&\text{ إذا وفقط إذا كان }
\ordeq{\tuple{\OLId{latin-looped}{A}, <}}{\tuple{\OLId{latin-looped}{B}_\OLId{latin-ordinary}{b}, \lessdot_\OLId{latin-ordinary}{b}}}
\text{ لبعض $b \in B$}
\end{align*}
وهذه التكافؤات من~\olref[iso]{ordisounique}،
و\olref[iso]{wellordinitialsegment}، و\olref[basic]{ordissetofsmallerord}.
\end{proof}

\end{document}
